\section*{Chapitre \ref{ch:b3:remote-sensing-gis} --- Télédétection et SIG}

\begin{solution}{exo:b3:remote-sensing-gis:1}
Champ~: $(0.42 - 0.06)/(0.42 + 0.06) = 0.75$, végétation verte dense, une culture
en pleine feuillaison. Champ labouré~: $(0.28 - 0.21)/(0.28 + 0.21) = 0.14$, presque
nu, un sol qui s'éclaircit légèrement vers l'infrarouge.
\end{solution}

\begin{solution}{exo:b3:remote-sensing-gis:2}
(a)~$185000/30 \approx 6170$ pixels~; $185000/15 \approx 12300$ dans la
bande panchromatique. (b)~$2^{12} = 4096$ niveaux~; $2^8 = 256$. (c)~16~jours
à lui seul~; avec Landsat~9 sur la même orbite, à 8~jours d'intervalle, 8~jours.
\end{solution}

\begin{solution}{exo:b3:remote-sensing-gis:3}
(a)~Une frange vaut $\lambda/2 = \qty{27.7}{mm}$~; six franges correspondent
à \qty{166}{mm} de variation selon la \omterm{def:b3:remote-sensing-gis:insar}{ligne de visée}, ici un éloignement par rapport au satellite, puisque
le bassin s'affaisse. (b)~$56.6/2 = \qty{28.3}{mm}$.
\end{solution}

\begin{solution}{exo:b3:remote-sensing-gis:4}
(a)~$t = 2d/c = 2\times\num{5e5}/\num{3.00e8} = \qty{3.34}{ms}$.
(b)~$c\,\delta t/2 = \num{3.00e8}\times\num{8e-10}/2 = \qty{0.12}{m}$~; une
surface échantillonnée par de nombreux photons est bien mieux localisée, la précision
décroissant comme l'inverse de la racine carrée de leur nombre.
\end{solution}

\begin{solution}{exo:b3:remote-sensing-gis:5}
(a)~En unités SI $L = \qty{9.0e6}{W.m^{-2}.sr^{-1}.m^{-1}}$. Avec $c_1 =
2hc^2 = \qty{1.191e-16}{W.m^2.sr^{-1}}$, $c_2 = \qty{0.014388}{m.K}$,
$\lambda^5 = (\num{1.09e-5})^5 = \num{1.539e-25}$, $c_1/(\lambda^5L) = 86.0$,
et $T_b = c_2/(\lambda\ln 87.0) = \qty{295.6}{K}$. (b)~La température cinétique
vérifie $0.98\,B_\lambda(T) = L$, c'est-à-dire $B_\lambda(T) = L/0.98$~: la même
formule donne $T = \qty{296.9}{K}$, soit \qty{1.3}{K} de plus.
\end{solution}

\begin{solution}{exo:b3:remote-sensing-gis:6}
(a)~$R = 693/\cos\ang{35} = \qty{846}{km}$. (b)~Antenne réelle~: $R\lambda/L =
\num{8.46e5}\times0.0555/12.3 = \qty{3.8}{km}$~; ouverture synthétique~: $L/2 =
\qty{6.2}{m}$, environ 600 fois plus fine. (c)~L'ouverture synthétique est la
longueur de trajectoire sur laquelle un point reste dans le faisceau, $R\lambda/L =
\qty{3.8}{km}$.
\end{solution}

\begin{solution}{exo:b3:remote-sensing-gis:7}
$\partial z/\partial x = (418 - 405)/60 = 0.217$ (est moins ouest),
$\partial z/\partial y = (418 - 401)/60 = 0.283$ (nord moins sud)~;
$|\nabla z| = 0.357$, soit une pente de $\arctan 0.357 = \ang{19.6}$. Le sol
monte vers le nord-est, il est donc orienté vers le sud-ouest~: $A =
\operatorname{atan2}(-0.217, -0.283) = \ang{217}$. Ombrage~: $\cos\ang{45}
\cos\ang{19.6} + \sin\ang{45}\sin\ang{19.6}\cos(\ang{315} - \ang{217}) =
0.666 - 0.033 = 0.63$, un peu plus sombre qu'un sol plat (0.71)~: la pente
tourne légèrement le dos à une source lumineuse située au nord-ouest.
\end{solution}

\begin{solution}{exo:b3:remote-sensing-gis:8}
Totaux par ligne 52, 48, 38, 22~; totaux par colonne 54, 46, 38, 22~; diagonale 133 sur
160~: exactitude globale \qty{83.1}{\percent}. Exactitudes de l'utilisateur~: eau
$48/52 = 0.92$, forêt $37/48 = 0.77$, cultures $29/38 = 0.76$, zone urbaine
$19/22 = 0.86$. Exactitudes du producteur~: eau $48/54 = 0.89$, forêt $37/46 = 0.80$,
cultures $29/38 = 0.76$, zone urbaine $19/22 = 0.86$. $p_e = (52\cdot54 + 48\cdot46 +
38\cdot38 + 22\cdot22)/160^2 = 0.271$, $\kappa = (0.831 - 0.271)/(1 -
0.271) = 0.77$. Les cultures sont le plus souvent omises~: 9 des 38 pixels réellement cultivés
ont été classés ailleurs, surtout en forêt.
\end{solution}

\begin{solution}{exo:b3:remote-sensing-gis:9}
(a)~$L = 0.04\,B_\lambda(1443) + 0.96\,B_\lambda(290)$ à
\qty{10.9}{\micro\metre}, \qty{28.6}{W.m^{-2}.sr^{-1}.\micro\metre^{-1}},
dont la \omterm{def:b3:remote-sensing-gis:brightness}{température de brillance} vaut \qty{396}{K}, à comparer avec une température moyenne de
$0.04\times1443 + 0.96\times290 = \qty{336}{K}$. (b)~À
\qty{2.2}{\micro\metre}, $B_\lambda(1443)/B_\lambda(290) = \num{6.8e7}$~:
la lave domine entièrement le pixel, et sa \omterm{def:b3:remote-sensing-gis:brightness}{température de brillance} est
\qty{845}{K}. La bande infrarouge à ondes courtes détecte les petits points chauds
bien mieux~; la bande thermique, sensible à l'ensemble du pixel, est celle qui
mesure le sol qui les entoure.
\end{solution}

\begin{solution}{exo:b3:remote-sensing-gis:10}
Avec le satellite à l'est, $d_1 = u_E\sin\theta + u_z\cos\theta$~; à
l'ouest, $d_2 = -u_E\sin\theta + u_z\cos\theta$. En additionnant et en soustrayant,
$u_z = (d_1 + d_2)/(2\cos\theta) = 102/1.638 = \qty{62.3}{mm}$ et $u_E =
(d_1 - d_2)/(2\sin\theta) = 26/1.147 = \qty{22.7}{mm}$. Un seul passage fournit une
équation pour deux inconnues~: tout mélange de soulèvement et de déplacement vers l'est
de même projection sur la \omterm{def:b3:remote-sensing-gis:insar}{ligne de visée} produit les mêmes franges.
\end{solution}

\begin{solution}{exo:b3:remote-sensing-gis:11}
(a)~D'après le \cref{thm:b3:remote-sensing-gis:fringe}, $\Delta\phi =
4\pi\Delta r/\lambda < \pi$ impose $\Delta r < \lambda/4 =
\qty{13.9}{mm}$. (b)~$u_z = k z_0(r^2 + z_0^2)^{-3/2}$ avec $k = (1 -
\nu)\Delta V/\pi$~; $\mathrm du_z/\mathrm dr = -3kz_0r(r^2 + z_0^2)^{-5/2}$,
dont la dérivée s'annule où $(r^2 + z_0^2) = 5r^2$, soit en $r = z_0/2$. On a alors
\[
\left|\frac{\mathrm du_z}{\mathrm dr}\right| = \frac{3k}{2z_0^3}\Big(\frac54\Big)^{-5/2}
= \frac{0.859\,k}{z_0^3} = \frac{0.859\times\num{2.39e6}}{\num{6.4e10}} = \num{3.2e-5},
\]
\qty{32}{mm} par kilomètre. (c)~$13.9/32 = \qty{0.43}{km}$~: des pixels de taille inférieure
à environ \qty{430}{m}, condition largement respectée par Sentinel-1.
\end{solution}

\begin{solution}{exo:b3:remote-sensing-gis:12}
(a)~Pour $N, R \ge 0$ non simultanément nuls, $|N - R| \le N + R$, donc $|\mathrm{NDVI}|
\le 1$~; remplacer $N$ et $R$ par $aN$ et $aR$ revient à multiplier le numérateur et
le dénominateur par $a$. (b)~Sans correction~: $(0.502 - 0.051)/(0.502 + 0.051) =
0.816$~; avec correction~: $(0.512 - 0.081)/(0.512 + 0.081) = 0.727$. (c)~La luminance de
trajet abaisse l'indice de la végétation verte, d'environ un dixième ici, et
d'une quantité qui varie avec la brume~: une carte d'images non corrigées
montrerait autant la météo que la végétation, et les comparaisons entre dates
seraient dénuées de sens.
\end{solution}

\begin{solution}{pb:b3:remote-sensing-gis:1}
\textbf{1.} Forêt $(0.38 - 0.04)/0.42 = 0.81$~; coulée de 1890 $0.23/0.37 =
0.62$~; de 1960 $0.10/0.32 = 0.31$~; de 2005 $0.02/0.26 = 0.08$.

\textbf{2.} 1890, puis 1960, puis 2005~: la coulée la plus ancienne est la plus verte,
dans l'ordre des âges~; elle n'a pas encore atteint la forêt.

\textbf{3.} La roche nue réfléchit un peu plus dans le proche infrarouge que dans le
rouge, comme le sable brun de la figure du chapitre~; les lichens et les premières
plantes ajoutent un peu~: 0.08 est la valeur d'un sol presque nu.

\textbf{4.} Une pente orientée à l'opposé de la lumière assombrit les deux bandes dans la même
proportion~; le rapport les annule, alors que la seule réflectance dans le proche infrarouge
confondrait une forêt ombragée avec une végétation clairsemée.

\textbf{5.} Par la \cref{prop:b3:remote-sensing-gis:planck-inversion} avec
$\lambda = \qty{10.9}{\micro\metre}$~: le plateau à \qty{291.1}{K}, le
pixel de la fumerolle à \qty{302.3}{K}.

\textbf{6.} $0.97\,B_\lambda(T) = \num{8.40e6}$ donne $T = \qty{293.0}{K}$.

\textbf{7.} $B_\lambda(\qty{370}{K}) =
\qty{22.48}{W.m^{-2}.sr^{-1}.\micro\metre^{-1}}$; $f = (9.95 - 8.40)/(22.48
- 8.40) = 0.110$, soit environ \qty{1100}{m^2} du pixel de \qty{e4}{m^2}.

\textbf{8.} Le rayonnement est une fonction convexe de la température
(\cref{prop:b3:remote-sensing-gis:mixed})~: la partie chaude ajoute plus de rayonnement
que sa part de température, si bien que le pixel paraît plus chaud que sa moyenne,
$0.110\times370 + 0.890\times291.1 = \qty{299.8}{K}$ contre
\qty{302.3}{K}.

\textbf{9.} La phase de chaque écho dépend du trajet exact~; un fort
changement de point de vue modifie la somme des échos des diffuseurs, et
l'\omterm{def:b3:remote-sensing-gis:insar}{interférogramme} perd sa cohérence~; il ajoute en outre une phase due au relief.
$48/12 = 4$ cycles pour un seul satellite.

\textbf{10.} $5.4\times27.73 = \qty{149.7}{mm}$ vers le satellite.

\textbf{11.} $u_z = 149.7/\cos\ang{35} = \qty{182.8}{mm}$.

\textbf{12.} Le sol se déplace vers l'extérieur en plus de monter~: du côté est
le mouvement vers l'extérieur va vers le satellite et s'ajoute au soulèvement, du côté ouest,
il se retranche (\cref{prop:b3:remote-sensing-gis:mogi-los}). La variation
tombe plus vite vers l'ouest, où les franges sont serrées~; le motif est étiré
vers l'est.

\textbf{13.} Sur cette ligne, le déplacement horizontal est nord--sud,
perpendiculaire à la visée, si bien que le radar ne voit que $u_z\cos\theta$, qui a
la forme du soulèvement. Le soulèvement se réduit de moitié en $r = 0.766\,z_0$, donc $z_0 =
3.4/0.766 = \qty{4.44}{km}$.

\textbf{14.} La moitié de 5.4~: 2.7 franges.

\textbf{15.} Une déformation est la même sur les deux \omterm{def:b3:remote-sensing-gis:insar}{interférogrammes} après
mise à l'échelle par les intervalles de temps, ou continue de croître de couple
en couple~; un motif atmosphérique change d'une date à l'autre et suit souvent
le relief. Un second couple montrant la même figure en cible confirme la
déformation.

\textbf{16.} $\Delta V = \pi u_z(0)z_0^2/(1 - \nu) = \pi\times0.1828\times
4436^2/0.75 = \qty{1.51e7}{m^3}$.

\textbf{17.} $\num{1.51e7}/48 = \qty{3.1e5}{m^3}$ par jour, \qty{1.1e8}{m^3}
par an si le rythme se maintenait.

\textbf{18.} $u_r \propto r/(r^2 + z_0^2)^{3/2}$ est maximal là où sa
dérivée s'annule, $r^2 + z_0^2 = 3r^2$, $r = z_0/\sqrt2 = \qty{3.14}{km}$~;
là, $u_r = \qty{70}{mm}$, soit \qty{38}{\percent} du soulèvement central.

\textbf{19.} Rayon \qty{1.48}{km}, volume $\frac43\pi(1.48)^3 =
\qty{13.6}{km^3}$~; la variation représente $\num{1.51e7}/\num{1.36e10} =
\qty{0.11}{\percent}$ de celui-ci.

\textbf{20.} Non~: un dixième de frange représente \qty{2}{\percent} du comptage, et
la profondeur, lue sur une demi-largeur de \qty{3.4}{km}, éventuellement jusqu'à
\qty{0.1}{km}, intervient au carré, soit \qty{6}{\percent}~; la correction de \qty{3}{\percent}
est comprise dans cette incertitude.

\textbf{21.} Le réchauffement du champ de fumerolles (partie~II) au fil du temps,
une hausse des petits séismes sous le renflement, une évolution des gaz
émis, et le soulèvement mesuré par positionnement par satellite ou par nivellement en
quelques points du flanc.

\textbf{22.} Le réservoir sous le mont Halvren a augmenté de
\textbf{15~millions~\unit[text-series-to-math=true]{m^3}} ($\num{1.51e7}\,\unit{m^3}$).
\end{solution}
